\chapter{Posplošeni linearni modeli}

Edini napovedni model, ki smo si ga ogledali do sedaj, je linearna regresija. Ponovimo: na vhodu imamo značilke z zveznimi vrednostmi $x_i$, ki jih utežimo z utežmi $\theta_i$ ter jim prištejemo prosti člen (intercept) $\theta_0$, s čimer dobimo napoved ciljne vrednosti $\hat{y}$. Gradnja napovednega modela poteka tako, da iščemo takšne vrednosti parametrov, ki minimizirajo vsoto kvadratov napak med dejanskimi in napovedanimi vrednostmi ciljne spremenljivke v učni množici. Linearna regresija je eden najenostavnejših regresijskih modelov; enostavnejši pristop bi bil le napovedovanje s povprečno vrednostjo odvisne spremenljivke v učni množici, vendar tak pristop ne upošteva vrednosti vhodnih značilk, zato ga težko obravnavamo kot pravi napovedni model

V tem poglavju nam bo koristil vektorski in matrični zapis zgornjega modela, kjer podatke zapišemo v matriko značilk $X$, parametre združimo v vektor $\boldsymbol{\theta}$, napovedi pa izrazimo kot $\hat{\mathbf{y}} = X \boldsymbol{\theta}$, kar omogoča kompaktnejši zapis in učinkovitejšo implementacijo algoritmov za učenje modela. Pri tem predpostavimo, da je prva kolona v matriki podatkov kar enotski vektor, da sicer ne kompliciramo s posebno obravnavo s prostim členom in da gredo indeksi za $\theta$ od $0$ do $d$, kjer je $d$ število neodvisnih spremenljivk.

Postavi se vprašanje, ali obstajajo drugi, podobno enostavni modeli, morda za nekoliko drugačne napovedne naloge, torej takšni, kjer bi značilke ravno tako povezali z uteženo vsoto, nato pa dobljeno vrednost preoblikovali z ustrezno (nelinearno) povezovalno funkcijo, tako da bi lahko modelirali tudi diskretne ali kako drugače omejene ciljne spremenljivke. V nadaljevanju začnemo s primerom takšnega modela, ki ga bomo uporabili za razvrščanje, nato pa se vprašamo, ali so tovrstne razširitve dovolj splošne za širši razred modelov, kaj pri njih pravzaprav predpostavimo, od kod izhajajo njihove kriterijske funkcije in ali gre pri tem za zanimiv razred modelov s skupnimi lastnostmi.

\section{Primer}

Začnimo z (izmišljenim) primerom. V tabeli~\ref{tab:bled} so zbrani kopalci na Bledu, ki smo jih vprašali, koliko ur na teden se ukvarjajo s športom in koliko ur so prejšnjo noč spali. Zabeležili smo tudi, če so v dnevu intervjuvanja uspeli odplavati na otok. Ta je od bližnjega kopališča oddaljen več kot pol kilometra v eni smeri in je plavanje na otok in nazaj kar zalogaj, ki ne bi bil ravno primeren za slabše kopalce. Cilj je razviti aplikacijo, ki bi kopalcem svetovala, sevedea glede na fizično pripravljenost in spočitost, ali naj se napotijo na tak podvig. Aplikacija seveda potrebuje napovedni model, tega pa lahko zgradimo iz naših pdatkov.

\begin{marginfigure}
\footnotesize
\centering
\begin{tabular}{lccc}
\toprule
Ime & Vadba & Spanje & Otok \\
\midrule
Alenka     & 7  & 8 & 1 \\
Ana        & 2  & 8 & 0 \\
Andrej     & 5  & 5 & 0 \\
Blaž       & 5  & 9 & 1 \\
Boštjan    & 7  & 5 & 0 \\
Goran      & 12 & 6 & 1 \\
Gregor     & 10 & 9 & 1 \\
Helena     & 4  & 9 & 1 \\
Irena      & 9  & 3 & 0 \\
Janez      & 5  & 6 & 0 \\
Jure       & 8  & 4 & 0 \\
Katarina   & 4  & 3 & 0 \\
Klara      & 3  & 9 & 1 \\
Luka       & 5  & 4 & 0 \\
Maja       & 9  & 6 & 0 \\
Marko      & 4  & 7 & 0 \\
Matej      & 4  & 8 & 1 \\
Miha       & 6  & 4 & 0 \\
Mojca      & 11 & 5 & 1 \\
Nika       & 2  & 5 & 0 \\
Nina       & 8  & 6 & 1 \\
Petra      & 3  & 6 & 0 \\
Polona     & 9  & 8 & 1 \\
Rok        & 10 & 6 & 1 \\
Sara       & 6  & 7 & 1 \\
Sašo       & 10 & 7 & 1 \\
Sebastjan  & 12 & 5 & 1 \\
Tatjana    & 12 & 9 & 1 \\
Tjaša      & 11 & 3 & 0 \\
Tomaz      & 12 & 5 & 1 \\
\bottomrule
\end{tabular}
\end{marginfigure}

Ker so podatki dvodimenzionalni, jih je najbolje izrisati v razsevnem diagramu. Označimo tudi razred. Že prvi pogled na izris kaže, da je morda moč dobre plavalce in te, ki se bolj kopajo, ločiti s črto, oziroma odločitveno mejo. Ta je linearna, in jo torej lahko zapišemo kot X theta = 0. Izris vključuje tudi tri nove obiskovalce, Sara, Martin, in Leon. Kateremu izmed njih bo naša aplikacija, oziroma model, svetovala, da lahko odplava do otoka in priplava nazaj?

% slika





\section{Nomogrami}
